\section{Five tensions in the specification}
\label{sec:synthesis}

Read against the evidence, the specification pulls in two directions at five points, and cheap typed judgement turns each from a philosophical question into an engineering default.
For each tension we give the most charitable reading of the text, the evidence that bears on it, and options, at least one of which preserves the text.
The choice among options belongs to the framework's authors and to those who would implement it.

\paragraph{T1. Understanding emotion in language versus detecting persons' emotional states.}
\Clause{4.3} states that it is ``entirely unnecessary'' for PAI to detect a person's emotional state, which is fleeting, personal, context-sensitive and touches on dignity; \clause{4.3.3.1} asks PAI to ``understand human emotional states and processes''; and \clause{4.3.3.3} illustrates semantic recognition with two questions about a little dog whose difference lies in the intent of the utterance.
Read together, the clauses admit a consistent interpretation: recognise the emotional meaning and intent of utterances, never infer a person's inner state.
The text does not state this boundary, however, and a typed model makes it almost free to attach an emotion score to every message, which is how the boundary would be crossed in practice.
The evidence on emotion-related judgement is thin and mixed: in one annotation study \jev{} was confident on most empathy items yet on its most confident items barely more accurate than on all items (0.383 against 0.371, with three balanced classes), a failure the authors found in a low-cost generative model as well~\cite{arx24574}, while on multilingual valence--arousal regression \jev{} with a supervised layer fitted to the labels matched the best systems~\cite{arx35293}.
The wider literature is clearer: inferring emotion from observable signals is scientifically weak~\cite{barrett2019emotional} and ethically fraught~\cite{stark2021ethics}.
\emph{Options:} state the boundary explicitly, so that a safeguard judges the semantics and intent of utterances but never assigns emotional states to persons; and reserve emotion-related judgement for PAI auditing its \emph{own} outputs, for example for the fabricated intimacy that \clause{4.3.2} names as manipulation.

\paragraph{T2. Binary verdicts versus a three-valued ethics.}
\pol{} is built on a love/hate polarity but adds the state of absence (\clause{4.3.2}) and warns that its concepts serve structural analysis, not moral labelling (\clause{1.5}).
A yes/no question such as ``is this hate language?'' discards both qualifications.
The evidence points the same way from several directions: option names can drive the verdict~\cite{arx26758,arx00346}; complementary probabilities do not sum to one where the model is unsure~\cite{arx33209}; an explicit ``I don't know'' option is rarely chosen when correct, and ``none'' only inconsistently (54\% of the time on a medical benchmark, 7\% when recognising it needed arithmetic)~\cite{arx34024,arx39496}; and a yes/no probability placed conflict and ignorance in the same band, whereas a Choice question naming both states separated them in four cases~\cite{zen22848952}\zmark{}.
It also shows what helps: a separate yes/no question asking whether the evidence suffices detected missing information better than confidence, although it was a weaker filter for accepting answers~\cite{arx01006}.
\emph{Options:} make every love/hate question three-valued, with absence or insufficient information as an explicit option, and ask sufficiency as its own question whose use rate is monitored; keep ``contested'' (people disagree) as a separate flag routed to people, distinct from ``absent''; and keep moral vocabulary out of option names by binding neutral identifiers to definitions and testing them by renaming.
These options implement \clause{1.5} and \clause{4.3.2} rather than change them.

\paragraph{T3. Full transparency versus closed models.}
\Art{4} and \art{5}(4) require that operating logic and decision processes be ``completely open and transparent'' and that technology ``should not be a black box''.
TypeSafe has published neither the weights, the architecture, the reward nor the training data of hosted \jev{}, and does not offer it for self-hosting~\cite{arx24574,arx28919,arx00346}.
A charitable reading is that \art{4} concerns the transparency of NaturalDAO's processes and resource flows, which logging a closed component's inputs, outputs, versions and calibration data partly satisfies; but such logs cannot show how the component reaches its outputs, and open weights alone would not provide the ``reasoning steps'' of \art{5}(5) either.
Open models make the readout inspectable, but in several comparisons they were less accurate or less robust: Laya showed the most severe name inversion and was near chance on interaction-risk screening~\cite{arx26758,arx33401}.
The gap is not fixed: larger open decoders were more robust to option-name swaps than \jev{}~\cite{arx00346}, Chinese models initialised from Laya slightly exceeded \jev{} on their authors' general benchmark and, after domain fine-tuning, in medicine, though they reached 92\% of \jev{}'s average accuracy across specialised domains~\cite{arx36965}, and an open RLCD recipe exists~\cite{arx38850}.
\emph{Options:} (a) use closed models only as external baselines and run consequential judgements on self-hostable models whose weights, versions and calibration data are pinned in each decision record; or (b) allow closed models as one of several independent detectors, with logged inputs, outputs and versions.
Option (a) trades measured accuracy and robustness for inspectability; that trade-off should be made explicitly and revisited as open models improve.
Publishing decision records also conflicts with keeping probabilities from attackers (\cref{sec:g2}), which \cref{sec:design} resolves by delayed and coarsened publication.

\paragraph{T4. AI-only adjudication versus correlated errors.}
\Art{2}(1) makes NaturalDAO autonomous ``on the premise of close collaboration with all humanity'', \art{10} gives disputes to AI fact-finding and ethical assessment subject to human questioning, and \art{9}(7) states that human supervision ``does not directly intervene'' in decisions, with issues raised by questioning handled by NaturalDAO itself.
The text does not specify how the premise of collaboration constrains individual decisions.
The evidence bears on the non-intervention principle in two ways.
Escalation can preserve errors: on rubric grading, low-cost generative evaluators repeated \jev{}'s most confident errors in 96\% of cases (an upper bound, from items selected as confident errors), and evaluators caught at most 5 of 130 confident misses in agent-security screening~\cite{arx29769,arx33401}.
But escalation can also work, when the next stage is a stronger reasoner or an independently built rule~\cite{arx26550,arx02048}, although a later study contests the first (\cref{sec:update}).
Whether a human end-point would do better on governance items is untested: no study compares human reviewers with typed models on such items, and the oversight literature shows that people also over-rely on confident systems and disagree among themselves~\cite{parasuraman2010complacency,green2022flaws,bansal2021whole}.
\emph{Options:} (a) keep non-intervention, but make questioning effective by publishing a standing random audit of high-confidence decisions, with a duty for NaturalDAO to respond publicly to each discrepancy and to escalate to independently built checkers rather than peers; or (b) use the amendment procedure of \art{11} to define a class of welfare decisions, such as withholding welfare (and, by the corresponding change to \clause{5.4.1}, blocking a person's speech), whose automated execution a human reviewer can suspend pending review.
Option (a) preserves the text but leaves consequential decisions with automated components, so it falls outside the detector role defined in \cref{sec:intro}; we list it because it preserves the text, not because the evidence supports it. Option (b) adds a check whose value is plausible from the due-process literature~\cite{citron2008technological} and close to existing law (below), but untested on this task; the reference design of \cref{sec:design} implements it.

\paragraph{T5. Assessing each person versus equal connection.}
Among the things governed AI ``can'' do, \clause{5.6} lists quantifying ``each person's Love Language and hate-language characteristics'' and guiding behaviour with incentives.
\Clause{4.3.1} holds that treating every person equally and treating every person's behaviour equally ``cannot be equated'', and the surrounding text refuses to label people by what they have done.
A typed model can score every utterance of every person for a fraction of a cent, so a per-person aggregate is technically trivial.
The same models answered value questions from a default cultural position~\cite{arx36399}, showed group bias in one informal forced-choice test~\cite{simmons2026saidno}\gmark{}, performed worse on low-resource languages (as typed and generative models alike did in one benchmark~\cite{arx37647}), and could be steered by a third party's opinion inserted into the content~\cite{arx31142}.
Aggregated over a person's history, such errors could fall unevenly on groups and become a durable label that others can manipulate; no study yet measures differential error by speaker group on behavioural judgements.
Experience with social scoring systems shows how such scores become infrastructure for control~\cite{creemers2018social,liang2018constructing}.
\emph{Options:} score events and behaviours, never persons, and keep no persistent per-person aggregates; or, if \clause{5.6} is implemented, require consent, access to one's own record, an appeal path, and published error rates by group and language for any per-person measure.

\paragraph{Regulatory parallels.}
None of these tensions is unique to \pol{}, and law in the EU already takes positions on several, directly or by analogy.
The EU AI Act~\cite{euaiact2024} prohibits emotion recognition in workplaces and educational institutions except for medical or safety reasons (Art.~5(1)(f)); because the Act defines an emotion recognition system as one that infers emotions from biometric data (Art.~3(39); Recital~44), inference from text alone, which is what a typed model performs, falls outside the prohibition, and the provision bears on T1 by analogy rather than application; and social scoring, meaning the evaluation of persons over time on their social behaviour or personal characteristics where the score leads to detrimental treatment in unrelated contexts or to treatment that is unjustified or disproportionate (Art.~5(1)(c)), bearing on T5.
For high-risk systems it requires effective human oversight (Art.~14), and it gives persons affected by decisions based on certain high-risk systems listed in its Annex~III, with legal or similarly significant effects, a right to an explanation of the system's role and of the main elements of the decision (Art.~86). These apply only to high-risk systems, and neither a speech filter nor a private organisation's welfare decisions is listed in Annex~III, whose point 5(a) covers decisions on public assistance taken by or for public authorities; point 8(a) does list AI used in alternative dispute resolution, which may reach the adjudication of disputes under \art{10} where its outcomes produce legal effects for the parties (Recital~61). So for the safety valve and welfare decisions they bear on T4 and on the explanation duty of \art{9}(2) by analogy, whereas the GDPR applies wherever personal data are processed.
The GDPR restricts decisions based solely on automated processing that produce legal or similarly significant effects and, where such decisions rest on a contract or explicit consent, requires at least the right to obtain human intervention, to express one's view and to contest the decision (Art.~22(1),~(3))~\cite{gdpr2016}, which is close to option (b) of T4.
For content moderation, the Digital Services Act requires hosting providers to give a statement of reasons for each restriction (Art.~17) and online platforms to run an internal complaint-handling system (Art.~20)~\cite{dsa2022}, and the Santa Clara Principles ask for notice, reasons and appeal~\cite{santaclara2021}; these are close analogues of a safety valve that blocks hate language.
The options above, no emotional-state inference about persons, no persistent per-person scores, human-suspendable blocks with reasons and appeal, and published decision records, would bring an implementation closer to these norms.

\Cref{fig:framework} shows which requirements each tension draws on.
